\section*{Experimental Design}

This section separates the organization of the empirical study from the general model and implementation details in Materials and Methods. The goal is to clarify how the final benchmark and targeted ablations are structured, which results are confirmatory five-fold evidence, and which analyses are exploratory only.

\subsection*{Overall benchmark design}

The active paper-facing experiment source is version \texttt{V3}. The final benchmark evaluates six datasets, nine methods, four message-passing operators, and five outer folds under one unified protocol. This yields the main comparison matrix used throughout the paper. For readability, the main text reports compressed summary tables and figures, while the complete matrix is moved to the appendix.

\subsection*{Biological core and supplementary benchmark split}

The benchmark is organized into two layers. The biological core consists of PROTEINS, DD, and ENZYMES, which anchor the main biomolecular claim of the paper. A supplementary robustness package extends the analysis to MUTAG, AIDS, and Mutagenicity. This split avoids mixing the strongest biological narrative with the broader molecular robustness checks and makes the final interpretation more disciplined.

\subsection*{Gate ablation design}

The first targeted ablation studies adaptive gate control on the strongest graph-level residual line. The target configuration is \texttt{AIDS + GraphResGNN + GINConv}, and the compared settings are a learnable gate and three fixed coefficients (\texttt{0.0}, \texttt{0.5}, and \texttt{1.0}). This experiment is designed to test whether adaptive residual injection is stronger or more stable than fixed-strength graph-level reuse.

\subsection*{Cross-gate ablation design}

The second ablation focuses on settings where a cross architecture is already locally favored. Six \texttt{dataset + model + operator} targets are evaluated under the same four gate settings: \texttt{AIDS + NodeCrossGNN + GATConv}, \texttt{DD + NodeCrossGNN + GATConv}, \texttt{DD + NodeCrossGNN + GINConv}, \texttt{ENZYMES + NodeCrossGNN + GATConv}, \texttt{MUTAG + GraphCrossGNN + GATConv}, and \texttt{Mutagenicity + NodeCrossGNN + GATConv}. The purpose is not to prove that learnable gating always wins, but to test whether adaptive gate control remains robust when cross exchange is already competitive.

\subsection*{Residual-routing ablation design}

The main mechanism study targets representative residual- and cross-oriented lines: \texttt{AIDS + GraphResGNN + GINConv}, \texttt{PROTEINS + GraphResGNN + SAGEConv}, \texttt{AIDS + NodeCrossGNN + GATConv}, and \texttt{DD + NodeCrossGNN + GATConv}. Five routing settings are compared under full five-fold evaluation: \texttt{learnable}, \texttt{topk\_0p25}, \texttt{topk\_0p5}, \texttt{sparse\_0p02}, and \texttt{sparse\_0p05}. These experiments are designed to answer whether mild filtering of residual injection can outperform dense learnable routing on selected targets.

\subsection*{Residual-parameter sweep design}

To extend the routing analysis, two additional settings are evaluated on the same four targets: \texttt{topk\_0p75} and \texttt{sparse\_0p1}. This sweep is used to test whether the optimal top-k ratio is target-dependent and whether more aggressive sparsity harms otherwise strong residual lines.

\subsection*{Exploratory sensitivity analysis}

Older single-fold sensitivity scans are retained as supporting analyses rather than as main evidence. These scans focus on cross models on PROTEINS, DD, and ENZYMES and vary depth, dropout, and learning rate on fold \texttt{0} only. In the final paper, they should be interpreted as optimization-trend evidence and clearly separated from the five-fold confirmatory experiments in \texttt{V3}.
